\section{Síntesis y ampliación}

\subsection{Conclusiones centrales}

\begin{importantbox}{Ocho conclusiones transferibles}
\begin{enumerate}
  \item AI coding es un ciclo cerrado de context, model, apply y validation; cada capa requiere un SLO distinto.
  \item El blast radius lo determina la runtime/state boundary, no el número de services.
  \item La Merkle-style sync, la content cache y la retrieval evaluation convierten el trabajo completo en delta work medible.
  \item Metadata, jobs, source objects y derived index requieren cada uno su propia source of truth y su propia rebuild path.
  \item Retry, cron, MVCC bloat y rebuild forman una retroalimentación positiva de incidentes; al recuperarse, primero se detiene al producer.
  \item El costo oculto de una migration es el re-embedding, la dual validation, el cache warming y la capacity temporal.
  \item El object storage desacopla la durabilidad del compute; la global inference sigue necesitando cache, region, provider eval y fallback.
  \item Security, abuse, tool approval e incident response son infraestructura para la continuidad del producto.
\end{enumerate}
\end{importantbox}

\subsection{Tarea práctica: diseñar un AI coding backend}

Completa el siguiente diseño para un sistema que soporte autocomplete, repo chat y multi-file agent:

\begin{enumerate}
  \item Define el SLO y el blast radius de indexing, inference, apply y data plane;
  \item Diseña la Merkle-style sync, la chunk/version key, la embedding cache y la deletion semantics;
  \item Enumera la source of truth, la rebuild path y la migration validation de metadata, job, source object y vector segment;
  \item Escribe el degraded mode para retry storm, database disk pressure, object-storage cache miss y provider rate limit;
  \item Establece self-hosted/frontier routing, offline/online eval, privacy, quota, tool approval e incident command.
\end{enumerate}

\begin{knowledgebox}{Criterios de aceptación}
Una propuesta de alta calidad no se limita a enumerar «vector DB + LLM + queue». Debe explicar el estado reconstruible, la amplificación por retry, la dual validation, la model version telemetry y cómo se comprenden y revocan los permisos sobre el código y la terminal.
\end{knowledgebox}

\subsection{Lecturas adicionales}

Se recomienda leer, en este orden, de Cursor: \href{https://cursor.com/blog/secure-codebase-indexing}{secure indexing}, \href{https://cursor.com/blog/instant-apply}{Fast Apply}, \href{https://cursor.com/blog/cursorbench}{CursorBench}, \href{https://cursor.com/docs/enterprise/privacy-and-data-governance}{privacy} y \href{https://cursor.com/docs/agent/security}{agent security}; después, de turbopuffer: \href{https://turbopuffer.com/blog/turbopuffer}{object-storage search} y \href{https://turbopuffer.com/blog/continuous-recall}{continuous recall}; por último, contrástalos con \href{https://www.postgresql.org/docs/current/routine-vacuuming.html}{vacuum} de PostgreSQL y el \href{https://docs.aws.amazon.com/AmazonS3/latest/userguide/Welcome.html}{object/consistency model} de Amazon S3.

\end{document}
